% Chapter 16: Production Deployment and High-Throughput Inference
% Part II of the book.
% STATUS: book-standard prose (Part II).

\chapter{Production Deployment and High-Throughput Inference}
\label{ch:part2-production-deployment-and-high-throughput-}
\label{sec:deployment}
% -----------------------------------------------------------------------------

A trained parameter vector is, in principle, an inexpensive object to apply;
making good on that principle is a separate piece of engineering from making the
network trainable. Everything before this chapter ran inside a training loop,
where the dominant cost is a gradient step and the forward pass is a subroutine
of it. Deployment inverts the accounting. The offline phase --- data generation,
the ensemble statistics and
coordinate convention of
\chapref{ch:part2-the-practitioner-s-paradigm-infinite-dimen}, the optimizer and
schedule of
\chapref{ch:part2-training-regimes-optimization-and-failure-}, the checkpoint ---
is paid once. The online phase is paid per query, at whatever rate the
application demands, and it must survive three things a training loop never
tests: leaving the interpreter that produced it, changing grid or batch size
without changing meaning, and answering, to a specified tolerance, the question a
solver would have answered. Once trained, neural operator
surrogates can be exported to optimized runtime environments for real-time
digital twins and aerodynamic sweeps; this chapter is about what that export
actually fixes and what it does not\index{model export}\index{digital twin}.

The Part~I result being operationalized is discretization invariance, stated as
Definition~\ref{def:ch06-discretization-invariance} and made exact for one
spectral layer by Proposition~\ref{prop:ch06-zeroshot}. Its content is that one
$\theta$ defines a map on functions rather than on arrays, so a model trained on
one grid applies on another. Deployment must carry that property through
serialization, and the warning of this chapter is that it does so automatically
only if the exported artifact still computes its sizes from the tensor it is
handed. The other two load-bearing imports are the per-layer
cost ledger \eqref{eq:ch8-costs} of
\chapref{ch:part1-comparative-theoretical-taxonomy-of-archit}, unfolded below into
concrete counts, and the three-term error split \eqref{eq:ch06-decomp}, which
governs what a published surrogate error number means once training is over.

Prerequisites: fluency with the PyTorch module API --- how a layer
reads its input's shape --- together with \chapref{ch:part2-fourier-neural-operators-fno-from-scratch}
and \chapref{ch:part2-canonical-pde-benchmarks-and-reference-wor}, and the device
vocabulary: a graphics processing unit (GPU) programmed through CUDA (Compute
Unified Device Architecture) executes a sequence of kernels launched from the host
at a small but nonzero fixed cost, and it holds tensor contents in device memory
that must be filled and drained around each call\index{CUDA (Compute Unified Device Architecture)}\index{graphics processing unit (GPU)}.
No vendor-specific toolchain
is assumed. The route is three sections: freezing the graph and exporting it
(Section~\ref{sec:ch27-export}), counting what one field costs and how batching
divides that cost (Section~\ref{sec:ch27-cost}), and reading the speedup table
against a stated accuracy budget (Section~\ref{sec:ch27-calibration}). The chapter
closes Part~II; Part~III takes up the reason any of this is worth doing.

% -----------------------------------------------------------------------------
\section{Freezing the Graph: Tracing, Export, and Runtime Specialization}
\label{sec:ch27-export}
% -----------------------------------------------------------------------------

A PyTorch model in training mode is an interpreter program: each call to
\texttt{forward} runs Python, which reads \texttt{x.shape}, loops over a module
list, and issues tensor operations to the device \cite{paszke2019pytorch}. A
deployed surrogate cannot carry an interpreter, a checkpoint file, and the
transitive closure of a scientific Python environment into the process that
needs it. \texttt{torch.jit.trace} resolves this by recording the sequence of
operations that a single execution performs and storing that sequence, together
with the weights, as a just-in-time (JIT) compiled graph that a C++ front end
loads and runs without Python\index{TorchScript}\index{just-in-time (JIT) compilation}. Tracing rather than source analysis is the right
tool here for a specific reason: the operator loop in the
\texttt{FNO2d} implementation iterates over the architecture, not over the data,
so the recorded operation sequence is the same for every input. Had any branch
depended on a tensor \emph{value} --- early stopping a rollout, selecting a mode
count from an energy threshold --- the trace would silently record whichever
branch the example input happened to take.

Listing~\ref{lst:ch27-torchscript} is the whole export, and the argument lives in
its first and last lines. The first is \texttt{model.eval()}. In the implementation of
\chapref{ch:part2-fourier-neural-operators-fno-from-scratch} this is a genuine
no-op: that model contains neither dropout nor batch normalization, so nothing
about what the forward pass computes changes when the flag is set. It is written
anyway because tracing records whatever the module does in the mode it is called
in, and a module that did own such layers would have them recorded in their
training behavior --- activations randomly zeroed on every query, running
statistics quietly updated by the traffic meant to read them --- so the archive
must be built from a model in inference mode rather than converted afterwards.
The decisive lines are the last three. The graph is traced on a
$(1, C_{\mathrm{in}}, 64, 64)$ example and then called on a
$(1, C_{\mathrm{in}}, 256, 256)$ input, and the printed shape is an assertion
rather than a comment. It must be read as a test, because the tracer's own
consistency check re-runs the \emph{example} input and therefore cannot report a
shape error at any other grid.

\begin{lstlisting}[language=Python, caption={TorchScript (JIT) export of the FNO-2D implementation: trace on a $64\times64$ example, save the archive, then evaluate on a $256\times256$ grid. The printed output shape is the assertion that the export contract demands --- it is only satisfied if the graph still derives its sizes from the live input.}, label={lst:ch27-torchscript}]
model.eval()
example_input = torch.randn(1, in_channels, 64, 64, device='cuda')

# Trace model computational graph
traced_fno = torch.jit.trace(model, example_input)
traced_fno.save("fno2d_surrogate_traced.pt")

# Verify zero-shot super-resolution on refined test grid in C++ / TorchScript
fine_input = torch.randn(1, in_channels, 256, 256, device='cuda')
out_fine = traced_fno(fine_input)
print("Super-resolution output shape:", out_fine.shape) # (1, out_channels, 256, 256)
\end{lstlisting}

Why that assertion can fail is worth understanding precisely, since it is the
mechanism by which a mathematically invariant model becomes a mesh-bound
artifact. Tracing freezes what Python computed, not what Python would compute.
In \texttt{SpectralConv2d}, the statement \texttt{B, C, H, W = x.shape} produces
four Python integers during the trace, and those integers are then consumed by
Python-level arithmetic: the size of the zero-filled output spectrum, the
half-width \texttt{W // 2 + 1}, the explicit \texttt{s=(H, W)} passed to the
inverse real transform, and, in the coordinate embedding of
\chapref{ch:part2-the-practitioner-s-paradigm-infinite-dimen}, the length of the
\texttt{linspace} that builds the grid channels. Each of those quantities enters
the archive as a constant fixed by the example input\index{shape specialization}.
The batch size is baked
the same way, which is why an example traced at $B=1$ cannot serve the batched
sweep that Section~\ref{sec:ch27-cost} argues for. Three remedies exist and they
differ in what they cost: derive the sizes inside the graph so that shape
computations become recorded operations rather than Python values; trace one
archive per served resolution, which is cheap because the weights are shared and
only the graph differs; or export through a graph description format such as ONNX
(Open Neural Network Exchange), in which the served dimensions are declared
dynamic on the way out rather than inferred on the way in. In each case the check in the
listing is what distinguishes a claim from a measurement.

Two further caveats belong next to the shape check. First, a correct printed
shape is a statement about the graph, not about accuracy. Proposition~\ref{prop:ch06-zeroshot}
guarantees equality of the spectral branch's output across grids only for inputs
band-limited inside the retained modes, and Figure~\ref{fig:super_resolution}
records the accompanying caveat: evaluating a coarse-trained model on a fine grid
is resolution transfer on a field the coarse grid already resolved, not recovery
of detail that the coarse grid never carried. A deployment that refines the query
grid expecting the surrogate to supply the missing content is asking the artifact
for something no theorem in this book grants. Second, what the archive carries is
weights and operations. The normalizer fields $\mu(x)$, $\sigma(x)$ of
\chapref{ch:part2-the-practitioner-s-paradigm-infinite-dimen} are buffers of data
rather than parameters of a layer, and if they are not explicitly registered or
serialized alongside the model, the deployed operator silently consumes
un-normalized inputs\index{normalization!state shipped with the artifact}. That failure is the dangerous kind: the output has the
right shape, the right spatial structure, and the wrong amplitude, and nothing in
the runtime reports it.

The failure modes named above --- a baked shape, an omitted normalizer --- are not
caught by a smoke test that only checks that the artifact runs, so it is worth
stating the contract the artifact must satisfy as an object an operations reviewer
can inspect. Implementation Pattern~\ref{imp:ch27-artifact} does that, in the order
in which its four components are checkable.

\begin{implementation}[The deployment artifact is a bundle, not a file]
\label{imp:ch27-artifact}
An archive that reproduces the trained operator's behavior consists of four
components, and each is a separate failure mode if omitted: (i)~the traced graph
and weights, traced at the shapes actually served; (ii)~the normalizer fields
$\mu$, $\sigma$ and the interpolation rule that carries them to a new grid;
(iii)~the coordinate convention, since the augmented channels of
\chapref{ch:part2-the-practitioner-s-paradigm-infinite-dimen} fix the domain's
orientation and scaling and a permuted grid channel is a different input
function; and (iv)~the acceptance metric with its threshold, stated in the
mesh-invariant relative $L^{2}$ form \eqref{eq:rel_l2_loss}. Taken together the
four say what the artifact is certified on; missing any one of them leaves a
model that runs and returns fields whose provenance nobody can state.
\end{implementation}

% -----------------------------------------------------------------------------
\section{What One Field Costs: Arithmetic, Memory, and Overhead}
\label{sec:ch27-cost}
% -----------------------------------------------------------------------------

Asymptotics answer ``how does it scale''; deployment asks ``how many
floating-point operations, how many bytes, how many milliseconds''. Write $N = HW$
for the number of grid unknowns, $d_v$ for the hidden channel width, $L$ for the
number of spectral layers, $m_x, m_y$ for the retained modes, and $c_{\mathrm{in}}$,
$c_p$, $c_{\mathrm{out}}$ for the input channels after coordinate augmentation, the
projection width, and the output channels. Counting a real two-dimensional fast
Fourier transform (FFT) at roughly $\tfrac{5}{2}N\log_2 N$ operations per channel, a
complex multiply at six real operations, and a multiply-accumulate at two, one forward
pass of the \texttt{FNO2d} of
\chapref{ch:part2-fourier-neural-operators-fno-from-scratch} costs
\begin{equation}
  \label{eq:ch27-cost}
  \begin{aligned}
    \Phi(N) \;=\;
    & \underbrace{5\,L\,d_v\,N\log_2 N}_{\text{forward and inverse FFTs}}
      \;+\; \underbrace{12\,L\,d_v^{2}\,m_x m_y}_{\text{mode multiplies}} \\
    & + \underbrace{2\,L\,d_v^{2}\,N}_{\text{local bypass}}
      \;+\; \underbrace{2\,N\big(c_{\mathrm{in}}d_v + d_v c_p + c_p c_{\mathrm{out}}\big)}_{\text{lift and projection}}
      \quad \text{operations},
  \end{aligned}
\end{equation}
per field. Equation \eqref{eq:ch27-cost} refines the FNO line of the ledger
\eqref{eq:ch8-costs}, which retains the first two groups: the two transform terms
are $\Theta(d_v N\log N)$ and the mode multiplication is $\Theta(K d_v^2)$ for
$K$ retained modes. The bypass and the pointwise lifting and projection are
$\Theta(d_v^2 N)$ and $\Theta(N)$, asymptotically dominated by the transforms, and
they are nonetheless worth writing out because at the resolutions a surrogate is
actually served at they are not small.

Take the reference configuration of that implementation, $d_v = 32$, $L = 4$,
with $m_x = m_y = 12$ and $c_p = 128$, on a $64\times64$ grid. The four groups of \eqref{eq:ch27-cost}
come to roughly $3\times10^{7}$, $7\times10^{6}$, $3\times10^{7}$ and
$4\times10^{7}$ operations, so one field costs of order $10^{8}$ operations, and
only about seven percent of them touch the learned spectral multipliers. On a
$256\times256$ grid the same count gives about $2\times10^{9}$: sixteen times the
unknowns, roughly seventeen times the work, the near-linearity in $N$ being exactly
what the $\log$ factor predicts\index{Fourier neural operator!inference cost}\index{einsum contraction}. Two structural facts follow. The second group of
\eqref{eq:ch27-cost} --- the Einstein-summation contraction (\texttt{einsum}) with
which \texttt{SpectralConv2d} rotates the retained modes, its coefficient of twelve
being six real operations for a complex multiply taken over the two multiplier
arrays of \eqref{eq:ch23-paramcount} --- contains no factor of
$N$, so refinement dilutes the very quantity that tensor factorization of the
multiplier arrays, examined in
\chapref{ch:part2-fourier-neural-operators-fno-from-scratch}, compresses; and the parameters, not the
arithmetic, are resolution-independent, which is the sense in which
\eqref{eq:ch8-costs} is a statement about a fixed $\theta$ rather than about a
deployed service.

Memory behaves less gracefully, because it is carried by activations rather than
by weights:
\begin{equation}
  \label{eq:ch27-memory}
  M(B) \;\approx\; B\,L_{\mathrm{live}}\Big[\,16\,d_v\,H\big(\lfloor W/2\rfloor+1\big)
  \;+\; 4\,d_v\,N\Big]
  \;+\; 16\,L\,d_v^{2}m_x m_y
  \quad \text{bytes},
\end{equation}
where $B$ is the batch size, the two bracketed terms are the pair of complex
spectrum buffers and the real hidden state of one layer, at eight bytes per
\texttt{torch.cfloat} entry for each of the two frequency arrays and four bytes per
single-precision real entry, $\lfloor W/2\rfloor+1$ is the half-spectrum width that a
real-input transform returns by conjugate symmetry, and $L_{\mathrm{live}}$ counts
layers whose working sets are simultaneously resident. The final term prices the
multipliers themselves, the two arrays per layer of $m_x m_y d_v^{2}$ complex
entries counted by \eqref{eq:ch23-paramcount}, again at eight bytes an entry.
At inference, with no autograd graph to keep every
layer's spectrum alive for a backward pass, $L_{\mathrm{live}}$ can be taken as
one; in training it is $L$. For the $64\times64$ configuration the per-field
working set is a little under two megabytes and the weight store is about nine; at
$256\times256$ the activation term grows sixteenfold to twenty-five megabytes while
the weight store does not move. As a cross-check, at $B=16$ with all four layers
live the frequency term of \eqref{eq:ch27-memory} reproduces the seventy
megabytes of frequency-domain intermediates that
\chapref{ch:part2-fourier-neural-operators-fno-from-scratch} counts directly for
\texttt{x\_ft} and \texttt{out\_ft} on the same configuration. Equation
\eqref{eq:ch27-memory} is why the batch ceiling is set by memory rather than by
arithmetic, and why a high-resolution sweep is issued in chunks rather than as one
call\index{memory footprint!of inference activations}.

The third term in the price is the one no operation count sees. Split the wall
time of a call into a part fixed by the call and a part proportional to the work:
\begin{equation}
  \label{eq:ch27-latency}
  T_{\mathrm{wall}}(B) \;=\; \underbrace{T_{\mathrm{fixed}}}_{\text{launches, host, input/output}}
  \;+\; B\cdot \underbrace{\frac{\Phi(N)}{\eta\,\Phi_{\mathrm{peak}}}}_{t_{\mathrm{field}}},
  \qquad
  \bar{t}(B) \;=\; \frac{T_{\mathrm{wall}}(B)}{B} \;=\; \frac{T_{\mathrm{fixed}}}{B} + t_{\mathrm{field}} .
\end{equation}
Here $\Phi_{\mathrm{peak}}$ is the device's peak rate and $\eta$ the fraction of
it that the realized kernels attain. Equation \eqref{eq:ch27-latency} separates
the two quantities that a benchmark table routinely conflates. \emph{Latency} is
$T_{\mathrm{wall}}(1)$, the time until one answer exists, which is what a
feedback controller or an intraoperative guide is subject to; \emph{throughput} is
$1/\bar{t}(B)$, the number of fields per second, which is what a design sweep or a
posterior sampling loop is subject to. They are optimized in opposite directions,
because $\bar{t}$ falls with $B$ while $T_{\mathrm{wall}}$ rises\index{latency and throughput}\index{kernel launch overhead}.

The arithmetic makes the fixed part visible. A device in the class named in
Table~\ref{tab:benchmarks_hardware} sustains single-precision work at tens of
teraflops per second, so the $1.2$ milliseconds reported for the Darcy row
corresponds to a realized throughput of order $10^{11}$ operations per second
against a per-field count of order $10^{8}$ --- well under one percent of peak.
That is the signature of a call bounded by kernel launches, host-side argument
handling, and memory traffic rather than by arithmetic: a forward pass of that
network issues a few tens of kernels, each with a launch and bookkeeping cost that
does not care how large the batch is. Two of the four groups of \eqref{eq:ch27-cost} reach
the device by different routes --- the pointwise convolutions through cuDNN (the
CUDA Deep Neural Network library), with kernels chosen by the framework's dispatch,
and the transforms through the platform's own FFT routine --- so a single $\eta$ in
\eqref{eq:ch27-latency} should be read as an effective value for the whole call
rather than as a property of any one operator\index{cuDNN (CUDA Deep Neural Network library)}.
The consequence is the batching rule\index{batching}. Since
$T_{\mathrm{fixed}}$ is divided across $B$ fields in \eqref{eq:ch27-latency}, the
per-field time falls almost proportionally until the arithmetic term overtakes it,
at a batch $B^{*} = T_{\mathrm{fixed}}/t_{\mathrm{field}}$ whose value depends on
$\eta$ and therefore on the shapes; for the configuration above it lands somewhere
between ten and a hundred fields. $B^{*}$ is a measurement, not an inference, and
the recipe is to bisect on batch size until $\bar{t}$ stops falling, then to stop.
Below $B^{*}$ the GPU is idle and the sweep is bounded by launches; above it the
queue depth merely trades response time for volume, and a latency-sensitive
consumer should stay at $B=1$ and pay the fixed cost knowingly.

Precision policy enters \eqref{eq:ch27-latency} at a different place than batching
does, and the caveats attached to automatic mixed precision (AMP) in
\chapref{ch:part2-training-regimes-optimization-and-failure-} transfer to
deployment unevenly rather than shrinking. Of the three situations named there as
lying outside what AMP can be configured around, only the third --- binary16
gradient range and the loss scaling that repairs it --- is a property of
differentiation, and it dies with the autograd graph. The other two are arithmetic
on the forward arrays and survive export untouched. The first, cancellation in a
finite-differenced physics-informed residual, applies to any artifact that reports a
residual at serving time, since subtracting nearly equal numbers stored with
$2^{-8}$ relative resolution gives a derivative whose error is set by the format and
not by the model. The second, the accumulation of the metric itself, becomes the
acceptance test rather than the training loss, but it is still an accumulation over
$B\,C_{\mathrm{out}}\,H\,W$ entries, and a rounding error no smaller than the unit
roundoff of a 16-bit format is of the same order as the percent-level error floors a
deployed surrogate is being asked to beat\index{automatic mixed precision (AMP)!at inference}.
Behind both sits the unscaled \texttt{rfft2}: a field of unit spatial mean
contributes a direct-current coefficient equal to the sum of every grid value, and
refinement raises every fixed-wavenumber coefficient with it, so one buffer must
hold that growing zero mode together with a band edge set only by the field's
spectral decay --- the two ends whose ratio decides which way a reduced format
breaks. That is the operation Listing~\ref{lst:ch27-torchscript} performs on
its last three lines. The policy of \chapref{ch:part2-training-regimes-optimization-and-failure-}
therefore holds at serving time unchanged --- the $1\times1$ lifting and projection in
\texttt{bfloat16}, the transform and the mode rotation of \eqref{eq:ch27-cost} in
single precision, and every reported error in single precision. Reduced precision is
consequently a lever on $\eta$ --- how much of the machine's peak rate the pointwise
groups actually deliver --- rather than on $T_{\mathrm{fixed}}$, and it changes
neither the mode counts nor the accuracy the artifact was certified at. As with
batching, the size of the gain is a measurement made on the deployed graph, not a
property of the architecture.

% -----------------------------------------------------------------------------
\section{Reading the Speedups, and Calibrating Surrogate Error Against Solver Budget}
\label{sec:ch27-calibration}
% -----------------------------------------------------------------------------

Table~\ref{tab:benchmarks_hardware} collects production inference timings across
four application classes on the NVIDIA A100\index{NVIDIA A100}. Each row
pairs one forward pass of a trained surrogate with one complete solve of the same
problem by a conventional method: a finite element method (FEM) assembly and solve
in FEniCS for the Darcy case\index{FEniCS}; a pseudo-spectral integration --- the method
whose Fourier multiplier machinery Part~I analyzed --- for the two-dimensional
turbulent Navier--Stokes case\index{pseudo-spectral method!as an inference baseline}; a
volume discretization of the compressible flow around a car
body, that is, computational fluid dynamics (CFD)\index{computational fluid dynamics (CFD)}, with the
geometry-informed neural operator (GINO)\index{Geometry-Informed Neural Operator (GINO)!deployment on a surface mesh}
of \chapref{ch:part2-graph-neural-operators-gno-and-irregular-g} on the surrogate side;
and a global atmospheric forecast at $0.25^{\circ}$ resolution against the
numerical integration of the European Centre for Medium-Range Weather Forecasts
(ECMWF) Integrated Forecasting System (IFS)\index{ECMWF Integrated Forecasting System (IFS)}
\cite{pathak2022fourcastnet}. The aerodynamic row alone reports the cited study's
cost--accuracy comparison instead of a paired wall-clock time, as the table
caption records.

\begin{table}[htbp]
\centering
\small
\caption{\textbf{Production inference benchmarks: Neural Operator vs. Classical Solvers.} Per-configuration wall-clock latencies as reported for the device named in the column heading; the final column is the quotient of the two preceding columns, except at the daggered aerodynamic row: $0.1$\,s per shape on a single NVIDIA V100 and the $2.6\times10^{4}$ figure are as reported by \cite{li2023geometry}, the latter a cost--accuracy ratio at $3\%$ drag error against a GPU volume-CFD baseline on the same device, whose wall-clock time the study does not pair with it.}
\label{tab:benchmarks_hardware}
\resizebox{\textwidth}{!}{\begin{tabular}{lcccc}
\toprule
\textbf{Application Domain} & \textbf{Classical Solver} & \textbf{Solver Latency} & \textbf{FNO Latency (A100)} & \textbf{Speedup Factor} \\
\midrule
2D Darcy Permeability Flow & Finite Element (FEniCS) & 1.45 s & 1.2 ms & $\approx 1,200\times$ \\
2D Navier-Stokes Turbulence & Pseudo-Spectral & 4.80 s & 2.1 ms & $\approx 2,280\times$ \\
3D Car Aerodynamics (GINO) & GPU Volume CFD & --- & $0.1$\,s (V100) & $2.6\times10^{4}$\,\dag \\
Global Weather Forecast (0.25$^\circ$) & ECMWF IFS Numerical & 120 min & 1.8 s & $\approx 4,000\times$ \\
\bottomrule
\end{tabular}}
\end{table}

What the reader should conclude from Table~\ref{tab:benchmarks_hardware} is not
the ranking of domains by multiplier. Each latency column spans more than three
orders of magnitude on its own --- the classical timings run from $1.45$ seconds
to two hours, the surrogate timings from $1.2$ milliseconds to $1.8$ seconds ---
while the multipliers stay inside
$1.2\times10^{3}$ to $2.6\times10^{4}$, a spread of a little over a factor of twenty: the
acceleration is far more stable than either of the costs that produce it, which is
the behavior a fixed-cost learned map exhibits against solvers whose price scales
with the accuracy demanded. Note also that the speedup column adds no information;
each entry is the quotient of the two latencies, so it can be checked and must be
reproduced from the timings rather than the reverse. Read as a set of four
quotients, the table says something a single benchmark cannot: the ratio is
generated mostly by how much work the classical column does per query, not by
anything in the surrogate's parameters\index{speedup factor}\index{benchmark!wall-clock comparison}.

That asymmetry is the reason a speedup is not an accuracy claim, and the reason
neither number licenses a decision on its own. With a conventional solver,
accuracy is purchasable at query time: refine the mesh, tighten the convergence
tolerance, take a smaller step, and the error in the delivered field falls, at a
rate set by the method's order and the solution's regularity --- the same
algebraic decay that \eqref{eq:ch06-rate} establishes for the surrogate's own
grid error, with the roles of the two sides reversed.
Inside a deployed surrogate there is no such dial.
The mode count $m_x, m_y$, the width $d_v$, and the depth $L$ are frozen in the
archive, so paying more compute per query buys no accuracy at all; accuracy is
bought offline, in data and retraining. Whatever the operator is wrong about, it is
wrong about at $1.2$ milliseconds and at twelve. A deployment therefore needs a
budget rather than a ratio, and the natural one is the triangle inequality applied
to the labeling data: if $u_{N_{\mathrm{lab}}}(a)$ is the discrete solve that
produced the training labels on grid $N_{\mathrm{lab}}$ and $u(a)$ the exact
solution, then
\begin{equation}
  \label{eq:ch27-budget}
  \underbrace{\|\mathcal{G}_\theta a - u(a)\|_{L^{2}}}_{\text{error against the truth}}
  \;\le\;
  \underbrace{\|\mathcal{G}_\theta a - u_{N_{\mathrm{lab}}}(a)\|_{L^{2}}}_{\text{error against the labels, reported in training}}
  \;+\;
  \underbrace{\|u_{N_{\mathrm{lab}}}(a) - u(a)\|_{L^{2}}}_{\text{discretization error of the labels}},
\end{equation}
with norms in $L^{2}(\Omega)$ and each term readable in the mesh-invariant form
\eqref{eq:rel_l2_loss}\index{error budget!surrogate against solver labels}.
Equation \eqref{eq:ch27-budget} is the deployment reading
of the decomposition \eqref{eq:ch06-decomp}: the reported validation error measures
the surrogate against a solver's output, not against the solution, and the second
term --- which the training loop never sees --- bounds how much of the first term
means. It follows that a surrogate can be certified only down to the accuracy of
its labels. Where a design decision is one the coarse solver would have made
identically, that is a feature and the surrogate is a cheaper way of asking the
same question; where the decision turns on the fine detail the labeling grid
discarded, no retraining on that grid reaches it, and the budget is spent on
labels.

The practical test inherits this structure. It is run on held-out inputs, in the
relative $L^{2}$ metric so that the number means the same thing on the grid used
for labeling and on the grid being served, and then repeated in the functional the
decision reads, because an $L^{2}$ tolerance says nothing about a drag coefficient,
a plume front position, or the anomaly correlation used to score a forecast\index{verification!of a deployed surrogate}. That
second step is where Part~III's fidelity columns come from, and where the residual
hypotheses of \chapref{ch:part1-open-theoretical-problems-and-mathematical}
still stand: every bound there is stated under a compactness hypothesis on the
support of the training measure (Assumption~\ref{ass:ch09-compact}), so nothing in
this book bounds what a surrogate does to an input
drawn from outside it, and an out-of-distribution
(OOD)\index{out-of-distribution (OOD) failure} field returns a plausible field rather than an error message. Export makes that
gap operationally sharper, not duller, since a runtime that cannot warn and a
model that cannot refuse are combined in the same artifact; the roadblocks of
\chapref{ch:part3-synthesis-matrix-and-real-world-deployment} begin exactly there.

\takeaway{Summary of Part II}{Part II provides the practical engineering foundation for neural operators: from tensor layout standards and continuous grid embeddings to complete PyTorch implementations of DeepONet, FNO, and GNO/GINO. By coupling spectral convolutions with stable training regimes and production runtime tracing, practitioners can deploy discretization-invariant surrogates that accelerate physical simulations by up to four orders of magnitude.}

% -----------------------------------------------------------------------------
% Bridge to Part III
% -----------------------------------------------------------------------------

Every cost in this chapter has been priced per field, with the number of fields
treated as given. That is the assumption a single application permits and the
assumption the rest of the book abandons. Part~III opens with
\chapref{ch:part3-the-multi-query-bottleneck-in-computationa}, which begins from
the algebraic system a classical solver must form for every query,
$\mathbf{A}(\mathbf{u};\mathbf{p})\mathbf{u} = \mathbf{f}(\mathbf{p})$ with $N$
between $10^{7}$ and $10^{9}$ unknowns --- two to five orders beyond the grids
\eqref{eq:ch27-cost} prices here, so the geometry, memory, and resolution
questions that Part~II handled on tori of a few thousand points return at a
different scale. That chapter's claim is about counts rather than recipes:
posterior sampling over a permeability or tissue-stiffness field, a Bayesian
inverse problem in the sense of \cite{stuart2010inverse}, is stated there to
consume $\mathcal{O}(10^{5}\text{--}10^{7})$ forward evaluations by Markov chain
Monte Carlo (MCMC), shape optimization explores thousands of geometries, ensemble
weather forecasting integrates fifty perturbed initial states twice daily, and
clinical guidance needs answers inside a procedure\index{Markov chain Monte Carlo (MCMC)}\index{multi-query surrogate problem!query counts at scale}.
Against those numbers the millisecond scale of
Table~\ref{tab:benchmarks_hardware} stops being a benchmark and becomes a budget,
and the pieces assembled in Part~II --- a batched, exported, calibrated surrogate
bundle --- become the objects that the aerodynamic, atmospheric, subsurface,
cardiac, and tomographic chapters of Part~III must each rebuild under their own
geometry and their own verification standard.

\vspace{1.5em}
